Robotic manipulation encompasses a broad range of interactions between a robot and its environment: grasping, pushing, poking, pivoting, and tossing objects; transporting liquid-filled containers; handling loads near or beyond rated capacity; and adapting to mechanical faults mid-task.
These interactions span a continuous spectrum of dynamic complexity.
At one end sits form-closure grasping, which is computationally attractive precisely because, once an object is secured, its state is predictable and control reduces to a geometry problem.
At the other end sits a class of interactions where robot and environment dynamics are inseparably coupled---where joint torques, object inertia, contact impulses, and fluid motion jointly determine task outcome.

Robotic manipulation research has invested heavily in the simple end of this spectrum, and for good reason.
Grasping is reliable, compositional, and well-studied.
But it is limited in scope---and the full range of tasks that manipulation actually demands requires reasoning about dynamics, not just geometry.
